\chapter{现代人工智能方法分层答案}

\begin{enumerate}[leftmargin=*]
  \item 对 tensor 先写语义 shape，例如 $(B,T,F)$；dtype 决定精度和混合精度规则；device 决定 CPU/GPU 内存位置；requires-grad 决定是否把运算放入 autograd 图。每次层间转换都断言 shape、dtype 和 device，否则“能运行”不等于对了时间/特征轴。
  \item 因果卷积只读取当前及过去窗口；RNN 通过递归状态压缩前缀；self-attention 对可见 token 做内容寻址；Transformer 是多头 attention、前馈、残差和归一化的堆叠。四者都必须显式传入 causal mask，否则模型可能用预测日后的 token 得到虚假收益。
  \item GNN 增加节点、边、邻接和消息传递顺序；DRL 增加动作、转移核、奖励、策略和长期价值。前者的时间快照决定可用关系，后者的环境反事实决定动作是否真正改变成交/库存；两者都不能只用静态监督损失理解。
  \item 两层 MLP 为 $h=\phi(W_1x+b_1)$、$\hat y=W_2h+b_2$、$L=\frac12(\hat y-y)^2$。链式法则给 $\nabla_{W_2}L=(\hat y-y)h^T$，$\nabla_hL=(\hat y-y)W_2^T$，$\nabla_{W_1}L=[\nabla_hL\odot\phi'(W_1x+b_1)]x^T$。残差层 $h+F(h)$ 的 Jacobian 为 $I+J_F$，保留恒等梯度路径但不保证完全消失被消除。
  \item 因果核 $(1,-1)$ 在约定“当前减前一期”下给 $(x_1-x_0,x_2-x_1,x_3-x_2)=(2,-1,3)$。对称卷积在时点 $t$ 读取 $x_{t+1}$，即使训练损失很低也违反实时信息集；mask 的边界和 padding 方式必须和生产推理一致。
  \item 无位置编码时，置换 token 会同步置换 $Q,K,V$，attention 输出只按同一置换重排，因此对序列顺序不敏感。两 token logits $(0,1)$ 的权重为 $(1/(1+e),e/(1+e))$；加入位置编码或 causal mask 后，权重才区分先后和可见性。
  \item GCN 一层为 $H'=\sigma(\tilde D^{-1/2}\tilde A\tilde D^{-1/2}HW)$；若构图使用预测日后才出现的边，泄漏发生在输入生成而不是网络层。Bellman 方程 $V(s)=\max_a\{r(s,a)+\gamma\mathbb E[V(s')\mid s,a]\}$ 还要求环境转移和奖励正确；固定历史价格回放不能提供未执行动作的 $s'$。
  \item 逐期构造序列后，先比较均值池化对原序列和反转序列的差异，再检查因果 CNN、RNN、attention 和 Transformer。padding mask 只屏蔽不存在的 token，causal mask 屏蔽未来 token；两者混淆会让模型在短序列上看起来正常而在真实长度上泄漏。

  \item 旧参数的梯度为 $(4,2,1,1)$。学习率 $0.1$ 后损失为 $0.605$，高于 $0.5$；学习率 $0.01$ 后输出 $1.781$，损失约 $0.3053$。过大步长跨过激活边界，不能保证下降。
  \item 每个决策日只加入此前已生效的边和节点属性。把静态全图与逐期图快照并行运行；若静态全图显著更好且包含未来公告边，这个差异就是泄漏反例，不是模型优势。
  \item 若模型先做时间均值再接线性层，它删除了顺序；没有位置编码、因果 mask、序列层和置换负例时，应称“随机特征/均值基线”。深度名称不能替代对信息路径的证明。
  \item 全量微调、任务头、LoRA 的参数量和适应能力不同；比较时冻结基础权重、tokenizer、语料截止和 prompt，报告训练数据是否含目标期之后的信息。LLM 回测尤其要检查权重记忆、检索语料发布日期和 prompt 中的未来标签。
  \item 构造两期、有限库存的小 MDP，枚举每个动作序列得到最优价值，作为 RL 解析或独立计算参照；再比较学习策略的动作和价值。固定价格回放只评价历史动作，不能回答另一订单会不会成交、改变库存或冲击价格。
  \item 开放项目的最低闭环是：合法数据截止、简单基线、时间切分、独立局部 解析参照、预测到仓位/动作、成本与容量、反例实验、对模型假设与结果的解释。模型架构只是其中一项，不是证据本身。
\end{enumerate}

\subsection*{两层网络的全部梯度：基础练习}
激活输入为 $-1$，$h=0$，故 $\partial_aL=\partial_bL=\partial_cL=0$；$\partial_dL=\hat y-y=-1$。零激活可使部分参数收不到梯度。

\subsection*{记忆衰减与注意力的数值含义：基础练习}
从 $h_0=0$ 得 $h_1=1,h_2=0.5,h_3=0.25$；第一期冲击在没有新输入时按几何速度衰减。

\subsection*{Bellman 递推与图上的重复平均：基础练习}
仍都为 5，因为 $A\mathbf1=\mathbf1$。常数向量是这种传播的不动点。
